\clearpage
\section{Benchmark Design Details}
\label{appendix:benchmark-design}

The benchmark contains 100 queries. Each item stores the question, category,
topic, expected documents, required citations, optional forbidden citations,
expected answer points, difficulty, and whether the item is expected to benefit
from graph or normative-hierarchy retrieval. Table~\ref{tab:appendix-benchmark-design}
explains the design role of each category.

\begingroup
\footnotesize
\setlength{\tabcolsep}{2.3pt}
\renewcommand{\arraystretch}{0.92}
\begin{longtable}{@{}L{0.19\textwidth}L{0.24\textwidth}L{0.27\textwidth}L{0.22\textwidth}@{}}
\caption{Benchmark category design.}
\label{tab:appendix-benchmark-design}\\
\toprule
Category & What it tests & Example question type & Expected system behavior \\
\midrule
\endfirsthead
\toprule
Category & What it tests & Example question type & Expected system behavior \\
\midrule
\endhead
\bottomrule
\endfoot
Definition QA & Legal definitions. & ``What is an employee under the Labor Code?'' & Retrieve the exact definition article or clause. \\
Direct QA & Single-provision legal lookup. & ``How much notice must an employee give before unilateral termination?'' & Return the direct provision and cite it. \\
Document guidance QA & Connection between primary law and implementing guidance. & ``Which circular guides labor-contract content?'' & Retrieve both the Labor Code basis and the lower-rank guidance. \\
Exception-based QA & Exceptions, prohibitions, or special cases. & ``When may an employee quit without notice?'' & Identify the exception rule and avoid unrelated general rules. \\
Procedure QA & Legal procedure, deadlines, and obligations. & ``What must the employer do after probation ends?'' & Retrieve the procedure rule and summarize the required action. \\
Scenario-based QA & Mapping a user fact pattern to legal provisions. & ``What happens if the employee unlawfully terminates the contract?'' & Select the legally relevant rule for the described situation. \\
Comparison QA & Comparing two legal regimes. & ``Compare severance allowance and job-loss allowance.'' & Retrieve evidence for both sides of the comparison. \\
Multi-hop QA & Questions requiring several linked provisions. & ``How is annual leave calculated for a partial year and does probation count?'' & Combine multiple provisions without dropping required citations. \\
Calculation or table lookup & Tables, appendices, and numeric lookup. & ``Female retirement age in 2026?'' & Retrieve article evidence plus the relevant appendix or table chunk. \\
Hard-negative citation QA & Avoiding close but wrong legal bases. & ``Holiday overtime pay, not overtime-hour limits.'' & Retrieve the required citation and keep forbidden citations out of the top results. \\
Paraphrased real-user QA & Informal user wording. & ``My company moved me to another job for almost three months; is that okay?'' & Map informal phrasing to the correct legal concept. \\
Labor dispute procedure QA & Labor dispute procedure and civil procedure links. & ``Can I sue over unlawful termination and which court basis applies?'' & Retrieve both labor-law and BLTTDS procedural evidence. \\
Out-of-corpus QA & Unsupported topics outside the indexed corpus. & ``What are the 2026 regional minimum wages?'' & Refuse or return insufficient context rather than guessing. \\
\end{longtable}
\endgroup

The category distribution is intentionally mixed. Easy definition and direct
questions test basic retrieval, while multi-hop, appendix, paraphrase,
hard-negative, and dispute-procedure questions stress the parts of the system
where legal QA is most likely to fail.
